\newpage
\section*{Abstract}
Alzheimer's disease is a progressive neurodegenerative condition that causes memory loss and cognitive decline, and early detection is essential to improve patient care and treatment planning. Structural MRI reliably captures the brain changes linked to Alzheimer's, yet many reported deep-learning systems inflate their accuracy by splitting the data at the level of individual 2D slices, which lets scans from the same patient appear in both training and testing. This work develops and benchmarks a family of fully three-dimensional models for Alzheimer's classification on the Alzheimer's Disease Neuroimaging Initiative (ADNI) dataset (1{,}132 T1-weighted MPRAGE volumes from 395 subjects), under a strict \emph{subject-level}, leakage-free evaluation protocol with five-fold cross-validation. Two clinically meaningful tasks are studied: the binary task of separating cognitively normal (CN) subjects from Alzheimer's (AD) subjects, and the harder three-class task of CN vs.\ mild cognitive impairment (MCI) vs.\ AD. Five architectures are compared on an identical pipeline and split: a 3D DenseNet-121 CNN, a MedicalNet-pretrained 3D ResNet, a 3D Vision Transformer (ViT), a Swin Transformer encoder, and a hybrid that fuses a 3D CNN with a 3D ViT to combine local anatomical detail with global contextual reasoning. The hybrid CNN--ViT gives the highest binary result at $82.2\% \pm 5.7$ accuracy (macro-F1 $0.80$), but it is statistically indistinguishable from the from-scratch DenseNet ($81.0\%$; paired fold-wise test $p \approx 0.5$), and both clearly exceed the pretrained ResNet ($69.8\%$) and the $65.4\%$ majority-class baseline. On the three-class task the best model (Swin, $57.7\%$) is only about five points above the $52.2\%$ majority-class baseline, with MCI the dominant source of confusion; the hybrid was not evaluated on this task, and a from-scratch ViT degenerated to never predicting the majority class --- a useful negative result about data-hungry transformers and stacked imbalance corrections on small medical datasets. Model decisions are examined with a suite of 3D explainability methods (Grad-CAM, Integrated Gradients, occlusion, LIME, and attention rollout); because the volumes are not skull-stripped and the methods disagree on individual subjects, these maps are reported as partial, qualitative evidence rather than proof of anatomically grounded reasoning. The contribution of this work is not a headline-chasing accuracy number but a rigorous, reproducible, leakage-free comparison of single versus hybrid 3D architectures for Alzheimer's classification, together with an explicit account of what the current evidence does and does not support.
